% Port aus SwissTeX 1.3.1 (Produktionsfork, test-stress.tex), unveraendert:
% keine \swisscover-/\swisstitle-Aufrufe, nichts an v2.0 anzupassen.
% Belastungstest fuer Listen, Fussnoten, Verweise und Randfaelle im Raster.
\documentclass{swisstex}
\setrunningtitle{SwissTeX \textbar{} Belastungstest}
\begin{document}

\swisstitle{Belastungstest}{Elemente im Raster}{Was bricht, was hält}

\section{Listen}

\subsection{Aufzählung}

Vor der Liste ein Absatz, damit die Grundlinie bekannt ist.

\begin{itemize}
\item Erster Punkt in einer Zeile.
\item Zweiter Punkt, der über zwei Zeilen läuft und deshalb den
      Zeilenumbruch innerhalb des Eintrags zeigt.
\item Dritter Punkt.
\end{itemize}

Nach der Liste ein Absatz.

\subsection{Nummerierung}

\begin{enumerate}
\item Erster Schritt.
\item Zweiter Schritt.
\end{enumerate}

Nach der Nummerierung ein Absatz.

\subsection{Verschachtelt}

\begin{itemize}
\item Aussen.
  \begin{itemize}
  \item Innen.
  \end{itemize}
\item Wieder aussen.
\end{itemize}

Nach der Verschachtelung ein Absatz.

\section{Fussnoten und Verweise}

\subsection{Fussnote}

Ein Satz mit einer Fussnote.\footnote{Der Text der Fussnote.} Danach folgt
weiterer Satz, damit sich die Wirkung auf die Grundlinien zeigt.

Ein zweiter Absatz nach der Fussnote.

\subsection{Dritte Gliederungsebene}

\subsubsection{Untergeordnet}

Text unterhalb einer dritten Ebene, sofern die Klasse sie kennt.

\section{Weitere Blöcke}

\subsection{Zitat}

\begin{quote}
Ein abgesetztes Zitat über zwei Zeilen, das prüft, ob die Umgebung das
Raster hält oder verschiebt.
\end{quote}

Nach dem Zitat ein Absatz.

\subsection{Beschreibung}

\begin{description}
\item[Erstens] Erklärung dazu.
\item[Zweitens] Weitere Erklärung.
\end{description}

Nach der Beschreibung ein Absatz.

\subsection{Vorformatiert}

\begin{verbatim}
zeile eins
zeile zwei
\end{verbatim}

Nach dem vorformatierten Block ein Absatz.

\subsection{Langer Absatz über den Seitenumbruch}

Dieser Absatz soll lang genug sein, um einen Seitenumbruch zu erzwingen und
damit zu prüfen, ob das Raster auf der Folgeseite wieder korrekt aufsetzt.
Der Inhalt selbst ist bedeutungslos und dient nur der Länge. Der Inhalt
selbst ist bedeutungslos und dient nur der Länge. Der Inhalt selbst ist
bedeutungslos und dient nur der Länge. Der Inhalt selbst ist bedeutungslos
und dient nur der Länge. Der Inhalt selbst ist bedeutungslos und dient nur
der Länge. Der Inhalt selbst ist bedeutungslos und dient nur der Länge. Der
Inhalt selbst ist bedeutungslos und dient nur der Länge. Der Inhalt selbst
ist bedeutungslos und dient nur der Länge. Der Inhalt selbst ist
bedeutungslos und dient nur der Länge. Der Inhalt selbst ist bedeutungslos
und dient nur der Länge. Der Inhalt selbst ist bedeutungslos und dient nur
der Länge. Der Inhalt selbst ist bedeutungslos und dient nur der Länge.

Ein Absatz nach dem langen Absatz.

\end{document}
